\documentclass[10pt,a4paper]{amsart}
\usepackage[margin=19mm]{geometry}
\usepackage{amsmath,amssymb,mathtools}
\usepackage[scr=rsfs,cal=euler]{mathalfa}
\usepackage{xfrac}
\usepackage[unicode,hidelinks,bookmarks=false]{hyperref}
\newcommand{\ie}{i.e.\ }
\newcommand{\cf}{cf.\ }
\newcommand{\resp}{respectively\ }
\newcommand{\Subsection}{\subsection}
\providecommand{\nobreakdash}{\nobreak\hbox{-}}
\setlength{\emergencystretch}{2em}
\multlinegap0pt
\usepackage{bm}
\usepackage{bbm}
\usepackage{dsfont}
\usepackage{xspace}
\usepackage{cancel}
\usepackage{mathtools}
\usepackage{amsthm}
\makeatletter
\@ifundefined{Theorem}{\newtheorem{Theorem}{Theorem}}{}
\@ifundefined{Lemma}{\newtheorem{Lemma}[Theorem]{Lemma}}{}
\makeatother
\def\R{\mathbb R}
\def\V{\mathcal V}
\def\e{\mathop{\mathrm{e}}\nolimits}
\def\ltwo{\mathop{\mathrm{L}^2}\nolimits}
\def\sech{\mathop{\mathrm{sech}}\nolimits}
\title[Topological Levinson's theorem and corrections at thresholds]{Topological Levinson's theorem and corrections at thresholds: the full picture in a quasi-1D example}
\author{T.T. Nguyen, D. Parra, S. Richard}
\date{Source: arXiv:2509.12684v1 (CC BY 4.0)}
\begin{document}
\maketitle

\noindent\emph{Source and licence.} Topological Levinson's theorem and corrections at thresholds: the full picture in a quasi-1D example. Version arXiv:2509.12684v1 (CC BY 4.0), distributed under the Creative Commons Attribution 4.0 International licence, as recorded in the arXiv licence metadata for this exact version and shown on the abstract page https://arxiv.org/abs/2509.12684v1. Full rights notice: licenses/arxiv-2509.12684-CC-BY-4.0.txt.\par\smallskip

\noindent\textit{Editorial scope.} Counted unit: the Lemma whose source label is \texttt{None} in the article (source0.tex lines 1528--1541, source label \texttt{None}), delivered together with the article's own complete original proof of it (source0.tex lines 1543--1636, one \texttt{proof} environment of 94 lines). No mathematics is written, repaired or re-derived: statement and proof are the article's own text line for line. Required context retained uncounted: none. Every definition, display and label of those lines is retained unchanged. Omitted from this package and not redistributed: 1--1527, 1542--1542, 1637--2957. Nothing inside a retained excerpt is omitted or altered: every retained body is delivered exactly as the article's lines are written, so no omission receipt of any kind accompanies this package. Citations: the retained excerpts carry no citation command at all, so no bibliography entry of the article is transcribed and no reference is redistributed. Reference adaptation: none was needed and none was made; every reference inside the delivered unit resolves inside it, so no unresolved reference or citation is produced. Printed slips kept literal: no printed word, symbol, hypothesis or number of the counted statement or of its proof was altered. Layout and class: the article's own class is \texttt{article} with packages including hyperref, amssymb, amsmath, amsthm, mathrsfs, venturis2, fontenc, comment, graphicx, tikz, tikz-3dplot, wasysym, color, amsfonts; the wrapper is the lane's pinned \texttt{amsart} editorial wrapper at 10pt on A4, which restates the theorem-like environments the retained text needs and adds the article's own macro definitions that the retained text needs (\texttt{\textbackslash R}, \texttt{\textbackslash V}, \texttt{\textbackslash e}, \texttt{\textbackslash ltwo}, \texttt{\textbackslash sech}), with extra wrapper packages (bm, bbm, dsfont, xspace, cancel, mathtools). The delivered numbering is the unit's own. Assets: no retained span contains a figure environment, a graphics-inclusion command, a PDF-page inclusion, a table or a TikZ picture; no image file is copied into this package, the package declares no asset dependency and no image file is redistributed with it. Licence: version arXiv:2509.12684v1 of this article is distributed under the Creative Commons Attribution 4.0 International licence, as recorded in the arXiv licence metadata for this exact version and shown on the abstract page https://arxiv.org/abs/2509.12684v1. A full rights notice is delivered at licenses/arxiv-2509.12684-CC-BY-4.0.txt. Characters: every retained excerpt is pure ASCII, so the delivered engine, XeLaTeX with its default Latin Modern text font, reports no missing character.
\begin{Lemma}
In $\ltwo(\R)$, the following equality holds for $s,t\in \R$
$$
\big[\V^0_N \Theta(\V^0_{N/2})^*\big](s,t)=
\frac{1}{\sqrt{2}\pi}\big(3+\tanh(s)\big)^{1/2}\;\!
\e^{(s+t)/2}\sech(s+t)\;\! \chi_+(t)\;\! b(t) + k(s,t),
$$
where $\chi_+\in C^\infty(\R)$ satisfies $\chi_+(t)=0$ for $t<-1$ and $\chi_+(t)=1$ for $t>1$, 
where
$$
b(t):=\big(\e^t+\e^{-t}\big)^{1/2}\big(\e^{t/2}+\e^{-t/2}\big)^{-1},
$$
and where $k(s,t)$ denotes the integral kernel of a Hilbert-Schmidt operator.
\end{Lemma}

\begin{proof}
By direct computations, one gets that
\begin{align*}
& \big[\V^0_N \Theta(\V^0_{N/2})^*\big](s,t)\\
& = -\frac{1}{\pi}\big(3+\tanh(s)\big)^{1/2} \left(\frac{1+\tanh(s)}{\cosh(s)}\right)^{1/2}
\frac{1}{-2+\tanh(t)-\tanh(s)}\left(\frac{1}{\cosh(t)}\right)^{1/2} \\
& = -\frac{1}{\pi}\big(3+\tanh(s)\big)^{1/2} \big(1+\tanh(s)\big)^{1/2}
\frac{\cosh(s)^{1/2}\cosh(t)^{1/2}}{-2\cosh(t)\cosh(s)+\sinh(t-s)} \\
& =  -\frac{1}{\pi}\big(3+\tanh(s)\big)^{1/2} \e^{s/2}
\frac{\e^{t/2}+\e^{-t/2}}{-2\cosh(t)\cosh(s)+\sinh(t-s)}\;\!\frac{1}{\sqrt{2}} \;\!b(t) \\
& =  \frac{1}{\sqrt{2}\pi}\big(3+\tanh(s)\big)^{1/2} 
\frac{\e^{(s+t)/2}+\e^{(s-t)/2}}{\cosh(s+t)+\e^{s-t}} \;\!b(t). 
\end{align*}
Since the multiplication operators defined by the functions
$s\mapsto \big(3+\tanh(s)\big)^{1/2}$ and $t\mapsto b(t)$
are bounded, the statement holds if one shows that the following kernels correspond
to Hilbert-Schmidt operators:
$$
R_1(s,t):=\frac{\e^{(s-t)/2}}{\cosh(s+t)+\e^{s-t}}
$$
and
$$
R_2(s,t):=\frac{\e^{(s+t)/2}}{\cosh(s+t)+\e^{s-t}} - \e^{(s+t)/2}\sech(s+t)\;\! \chi_+(t).
$$

In order to compute the $\ltwo$-norm of $R_1$ and of $R_2$, one looks at these expressions in polar coordinates. Namely, for $R_1$ and for $r>0$ and $\omega \in [0,2\pi]$ one considers  the expression
$$
R_1\big(r\sin(\omega),r\cos(\omega)\big)=
\frac{1}{\e^{r(\cos(\omega)-\sin(\omega))/2}\cosh(r(\cos(\omega)+\sin(\omega)))+\e^{-r(\cos(\omega)-\sin(\omega))/2}}.
$$
By considering the equalities $\cos(\alpha)\pm\sin(\alpha)=\sqrt{2}\cos(\alpha\mp\pi/4)$
one infers that
the denominator of the previous expression satisfies
\begin{align*}
& \e^{r(\cos(\omega)-\sin(\omega))/2}\cosh(r(\cos(\omega)+\sin(\omega)))+\e^{-r(\cos(\omega)-\sin(\omega))/2} \\
& =\e^{\sqrt{2}r\cos(\omega+\pi/4)/2}\cosh\big(\sqrt{2}r\cos(\omega-\pi/4)\big)+\e^{-\sqrt{2}r\cos(\omega+\pi/4)/2}.
\end{align*}
Then, one observe that one of the two terms $\e^{\pm\sqrt{2}r\cos(\omega+\pi/4)/2}$ is exponentially growing (as a function of $r$) with a constant independent of $\omega$, except in a neighborhood of $\omega=\pi/4$ and of $\omega=5\pi/4$.
However, in these neighborhoods, the factor $\cosh\big(\sqrt{2}r\cos(\omega-\pi/4)\big)$ is exponentially growing. Thus, there exist $c>0$ and $d>0$ independent of $\omega$ such that
$$
R_1\big(r\sin(\omega),r\cos(\omega)\big)\leq c\e^{-dr}
$$
for all $\omega \in [0,2\pi]$. 

For $R_2$, observe that
\begin{align*}
R_2(s,t) & = \frac{(1-\chi_+(t))\cosh(s+t)-\chi_+(t)\e^{s-t}}{\e^{-(s+t)/2}\cosh(s+t)[\cosh(s+t)+\e^{s-t}]} \\
& = \frac{1-\chi_+(t)}{\e^{-(s+t)/2}[\cosh(s+t)+\e^{s-t}]}-\frac{\chi_+(t)\e^{s-t}}{\e^{-(s+t)/2}\cosh(s+t)[\cosh(s+t)+\e^{s-t}]}.
\end{align*}
For the first term and by using the same approach, observe that its denominator
can be rewritten as
\begin{align*}
&\e^{-(s+t)/2}[\cosh(s+t)+\e^{s-t}] \\
&=\e^{-\sqrt{2}r\cos(\omega-\pi/4)/2}
\big[\cosh\big(\sqrt{2}r \cos(\omega-\pi/4)\big)+\e^{-\sqrt{2}r\cos(\omega+\pi/4)}\big].
\end{align*}
In addition, because of the cut-off $1-\chi_+$ we can restrict our attention
to $\omega \in [\pi/2-\varepsilon,3\pi/2+\varepsilon]$ for some fixed $\varepsilon>0$
small enough.  In this region, the product 
$$
\e^{-\sqrt{2}r\cos(\omega-\pi/4)/2} \cosh\big(\sqrt{2}r \cos(\omega-\pi/4)\big)
$$ 
is exponentially growing (as a function of $r$) with a constant independent of $\omega$, except in a neighborhood of $\omega=3\pi/4$.
However, in any such neighborhood, the product
$$
\e^{-\sqrt{2}r\cos(\omega-\pi/4)/2}\e^{-\sqrt{2}r\cos(\omega+\pi/4)}
$$
is exponentially growing (as a function of $r$) with a constant independent of $\omega$.
One then infers that the first term is $R_2$ is smaller than $c\e^{-dr}$ for some
$c,d>0$ and all $\omega \in [\pi/2-\varepsilon,3\pi/2+\varepsilon]$.

The second term in $R_2$ is treated similarly, for $\omega$ restricted to $[\pi/2+\varepsilon,3\pi/2-\varepsilon]$ with  $\varepsilon>0$
small enough.  We consider
\begin{equation}\label{eq_HS1}
\frac{\e^{-\sqrt{2}r\cos(\omega+\pi/4)}}{\e^{-\sqrt{2}r\cos(\omega-\pi/4)/2}\cosh(\sqrt{2}r\cos(\omega-\pi/4))[\cosh(\sqrt{2}r\cos(\omega-\pi/4))+\e^{-\sqrt{2}r\cos(\omega+\pi/4)}]}.
\end{equation}
Since
$$
0<\frac{\e^{-\sqrt{2}r\cos(\omega+\pi/4)}}{\cosh(\sqrt{2}r\cos(\omega-\pi/4))+\e^{-\sqrt{2}r\cos(\omega+\pi/4)}}\leq 1
$$
the exponential decay of \eqref{eq_HS1} (as a function of $r$) with a constant independent of $\omega$ is obtained by the first two factors of the denominator, except in a neighborhood of $\omega=7\pi/4$. In any such neighborhood, observes that
$$
\cosh(\sqrt{2}r\cos(\omega-\pi/4))\big[\cosh\big(\sqrt{2}r\cos(\omega-\pi/4)\big)+\e^{-\sqrt{2}r\cos(\omega+\pi/4)}\big]\geq 1
$$
while 
$$
\frac{\e^{-\sqrt{2}r\cos(\omega+\pi/4)}}{\e^{-\sqrt{2}r\cos(\omega-\pi/4)/2}}
= \e^{-r\big(\sqrt{2}\cos(\omega+\pi/4)-\sqrt{2}\cos(\omega-\pi/4)/2\big)}\leq \e^{-dr}
$$
for some $d>0$ and independent of $\omega$.

By collecting the various estimates, one obtains that the functions $R_1$ and $R_2$ belong 
to $\ltwo(\R^2)$, and therefore the corresponding kernel $k$ in the statement is a Hilbert-Schmidt operator.
\end{proof}

\end{document}
